\chapter{Literature Review}
\label{sec:background}

This chapter positions the thesis at the intersection of large language
models, prompt decomposition, vehicle telematics analytics, and graph-based
evaluation. Telemetry governance remains relevant because vehicle data can be
sensitive, but the scientific focus is the evaluation of local open-source
LLMs as graph proposal systems for directed relationship recovery.

\summarynote{Literature role}{The chapter motivates a focused graph-recovery
study: LLMs propose directed relationships between telemetry features, prompt
decomposition is treated as an empirical design variable, and every proposed
graph is evaluated against an explicit ground-truth adjacency matrix.}

\section{LLMs as Analytical Interfaces}

Large language models can adapt to tasks from natural-language context and
structured instructions \parencite{brown2020language}. This makes them
attractive as analytical interfaces for operational data: instead of requiring
a user to inspect tables, scripts, or dashboards, an LLM can translate a task
description and a compact data summary into a structured output. In this
thesis, that output is not a prose summary but a directed edge list over
telemetry features.

This use creates a methodological tension. The model can combine feature
semantics with broad domain knowledge, but its output is not a validated
statistical estimate. A plausible edge may be wrong, reversed, transitive, or
merely correlational. LLM output must therefore be treated as a hypothesis to
be parsed and evaluated, not as an authoritative causal discovery result.

\section{Vehicle Telematics Analytics}

Vehicle telemetry provides structured measurements about energy state,
movement, acceleration, cumulative usage, and distance. For fleet monitoring,
these signals are valuable because operational decisions often depend on
relationships between variables rather than isolated readings. Acceleration,
speed, distance, mileage, and remaining electric energy form a compact system
in which feature interactions matter for diagnostics and monitoring.

This motivates graph-based representation. A directed graph can encode a
candidate relationship such as acceleration influencing speed, speed
influencing travelled distance, or driving behaviour influencing remaining
energy. Such graphs can support inspection, knowledge-graph export, and
downstream decision support. However, the graph remains an evaluated artifact:
in this thesis it is compared to a predefined Tesla EV ground truth, and no
claim is made that the LLM independently proves real-world causality.

\section{Prompt Engineering and Prompt Decomposition}

Prompt engineering shapes model behaviour by changing the input structure
without changing model weights. Direct prompting states the task and output
format. Structured prompting adds constraints, policies, decomposition, or
answer templates. Chain-of-thought prompting demonstrated that prompting can
change reasoning behaviour \parencite{wei2022chain}, while least-to-most
prompting frames complex tasks as ordered sub-problems
\parencite{zhou2023least}. These methods motivate the thesis design, although
the thesis evaluates only final parsed outputs, not hidden reasoning traces.

The evaluated prompt decomposition strategies are least-to-most prompting and
skeleton-of-thought prompting. Least-to-most prompting asks the model to move
from simpler physical relationships to more difficult accumulated-state
relationships. Skeleton-of-thought prompting groups the graph into
energy-state edges, driving-behaviour edges, cross-group edges, and an
integrity gate. Both are compared with zero-shot, no-instruction, and
policy-first baselines.

The relevant research question is not whether these prompts make the model
``reason'' in a human sense. The observable question is whether they improve
the predicted adjacency matrix under fixed data, model, parser, and metric
conditions.

\section{Structured Output and Data-to-Text Context}

The task also relates to structured-data generation. Natural language
generation research studies how non-linguistic inputs are selected,
organized, and verbalized \parencite{gatt2018survey,wiseman2017challenges}.
This thesis uses the same concern for content selection but changes the target
from text to graph structure. The model must select which feature
relationships to output and express them in a strict machine-parseable format.

This output-format constraint is central. The evaluator only counts edges that
can be parsed into the allowed feature vocabulary. If the model produces
prose, unsupported variables, or ambiguous relations, the parser may discard
them. The evaluation therefore measures a combined capability: relationship
proposal plus compliance with the edge-list contract.

\section{Graph-Recovery Evaluation}

Graph recovery can be evaluated by comparing a predicted adjacency matrix with
a ground-truth adjacency matrix. In this thesis, each model-strategy cell
produces a binary matrix over six Tesla EV features. The evaluation then uses
precision, recall, F1 score, structural Hamming distance, false discovery
rate, Pearson correlation, Hamming distance, and Frobenius norm.

This metric suite is needed because graph recovery has multiple failure
modes. A model can underpredict and miss true edges. It can overpredict and
produce false positives. It can reverse directions, which is especially
important in telematics because speed-to-acceleration and
acceleration-to-speed express different structural assumptions. It can also
output no parseable edges, yielding deceptively low false-discovery rate but
zero useful recall.

\section{Self-Consistency as Related but Unevaluated Work}

Self-consistency methods typically sample multiple model outputs and aggregate
them, for example by majority vote over answers or edge sets. The supervisor
notebook contains repeated-run and self-consistency-style machinery, but the
validated primary result JSON files used in this thesis report
\texttt{n\_runs = 1}. Therefore, self-consistency is treated as related future
work, not as an evaluated contribution. This boundary prevents the thesis from
claiming variance estimates or repeated-run improvements that are not present
in the supplied artifacts.

\section{Operational Governance Context}

Vehicle telemetry can expose sensitive operational patterns. This is why
local inference, data minimization, and artifact retention remain relevant to
deployment practice. Risk taxonomies for language models emphasize that
deployment context shapes downstream harms
\parencite{weidinger2022taxonomy,bender2021dangers}. For the present thesis,
these issues form a governance boundary rather than a result metric:
telematics graph-recovery systems should minimize unnecessary data exposure,
preserve artifacts for audit, and avoid treating LLM-generated graphs as
validated physical truth.

\section{Positioning of This Thesis}

Table~\ref{tab:related-positioning} summarizes the thesis position relative
to adjacent research areas.

\begin{table}[!htbp]
  \centering
  \small
  \setlength{\tabcolsep}{3pt}
  \begin{tabularx}{\textwidth}{@{}L{35mm}Y Y@{}}
    \toprule
    Area & Typical emphasis & Position of this thesis \\
    \midrule
    LLM prompting & Instruction following, reasoning support, decomposition & Evaluates prompt strategies by their parsed graph output, not by hidden reasoning. \\
    Vehicle telematics & Monitoring, diagnostics, energy and motion signals & Uses a Tesla EV feature set as a controlled graph-recovery task. \\
    Structured-data generation & Content selection and faithful output from structured inputs & Requires strict edge-list output over an allowed feature vocabulary. \\
    Graph evaluation & Predicted structure versus reference structure & Reports adjacency-matrix metrics for every model-strategy cell. \\
    Telemetry governance & Risk management and data minimization & Kept as operational context, not the central result metric. \\
    \bottomrule
  \end{tabularx}
  \caption{Positioning of the thesis relative to adjacent research areas. The contribution is a controlled evaluation of prompt decomposition for telematics graph recovery using open-source LLMs.}
  \label{tab:related-positioning}
\end{table}

\section{Research Gap}

Prior work motivates prompting, decomposition, and structured generation, but
it does not by itself answer whether a local LLM can recover a directed
telemetry graph when its output is forced into a parseable edge list and
judged against a fixed reference DAG. The gap addressed here is therefore
narrow but concrete: prompt strategies are evaluated as interventions in a
prompt-to-graph pipeline, not as general improvements in reasoning quality.

The thesis fills this gap with a controlled Tesla EV evaluation: two local
open-source models, five prompt strategies, one six-feature ground-truth DAG,
and eight graph-recovery metrics. The next chapters convert this positioning
into the task definition, methodology, and results.
